% SPDX-License-Identifier: Apache-2.0
\section{Related Work}
\label{sec:e-related}

Three systems bound this design, and it is closer to each than to the
others.

\textbf{Chisel}~\cite{bachrach2012} is the structural half. A hardware
module is a class in a general-purpose language, a bundle of wires has a
direction per field through \code{Flipped}, control flow on a signal is
\code{when} and \code{Mux}, and elaboration is the program running and
recording a graph. Sections~\ref{sec:e-units} to~\ref{sec:e-binding} are
that architecture in Rust, and they should be read as a port rather than
a proposal. SpinalHDL follows the same shape in the same host language.
Two things differ, both because of the host. Rust's ownership makes a
second driver on a wire a move error at compile time, where Chisel
detects it at elaboration. And Rust's trait system lets a configuration be
a type, so a build is a type alias and a binding path cannot go stale.

\textbf{Calyx}~\cite{calyx} is the nearest ancestor of the behavioural
half, and it is an intermediate representation rather than a language.
It separates structure, the wires and groups that compute, from control,
a program that says when groups run. Its control has \code{seq}, which
runs groups one after another, and \code{par}, which runs them at once.
Section~\ref{sec:e-units} arrives at the same two operators from the other
direction: an \code{await} is a \code{seq} boundary, and \code{join!} is
\code{par}. Calyx lowers both to a state machine driving the groups, and
that lowering is the one this paper has not yet done.

\textbf{Filament}~\cite{filament} types each signal with the interval of
cycles over which it is valid, so that composing two components with
different latencies is checked rather than simulated. That is what an
\code{async} operator whose latency the mapping decides is reaching for.
Filament makes the latency part of the type and checks it statically;
this paper defers it to the mapping and would learn it at elaboration.
Filament's is the stronger guarantee, and the price is that latency has
to be stated in the type. Whether an \code{async fn} can be given a
timeline type after the fact is an open question.

\textbf{RHDL}~\cite{rhdl} is the nearest thing in Rust, and the source
of one decision here. It embeds a hardware description language in Rust
with the clock domain as a type parameter of the signal, so that a
crossing without a synchroniser is a type error. Section~\ref{sec:e-wiring}
takes that directly, and departs from it only in what the default domain
means. RHDL's structural model is Chisel's; its behavioural model is
Rust functions on values, evaluated to produce combinational logic, with
registers explicit. It has no \code{async}, so the sequencing half of
this paper is as far from RHDL as from Chisel.

\textbf{XLS}~\cite{xls} is the one to weigh most carefully, because it
is a working toolchain that has done what Section~\ref{sec:e-next} calls
the next experiment. Its front end, DSLX, is a dataflow language with
Rust's surface syntax: \code{fn}, \code{let}, \code{match}, structs,
enums and parametric types, read by a parser of its own. A DSLX function
is scheduled into a pipeline by the compiler, with the stage boundaries
chosen to meet a clock period rather than placed by the designer, which
is what Section~\ref{sec:e-pipeline} asks of an \code{async fn}. A DSLX
\code{proc} is a process with state and channels, sending and receiving
in a loop, which is the sequenced half of this paper under another
name. XLS emits Verilog, and it has been used in silicon.

That makes XLS the nearest ancestor of the behavioural half, nearer than
Calyx, and it also makes it the counterexample to this paper. XLS is the
path the project's second thesis~\cite{theses} chose: a new language,
shaped like Rust, with its own toolchain. The merged specification of
the companion paper~\cite{merge} is on that path. This paper leaves it,
and XLS is the standing evidence that the path it leaves works. The
argument for leaving it is therefore not that a separate language cannot
be made to lower; XLS shows it can. It is that the type system, the
module system, the generics, the package manager and the editor arrive
already written when the host is Rust itself, and that a DSLX-shaped
language has to rebuild each of them. Whether that trade is worth the
one wall of Section~\ref{sec:e-walls} is the question the next
experiment answers.

What is not near. Bluespec~\cite{nikhil2004} schedules guarded atomic
rules, which is a different model of concurrency from processes that
await. High-level synthesis from C or C++ starts from sequential code and
infers a pipeline, which is the same ambition as
Section~\ref{sec:e-pipeline}, from a language that has no way to say
where a boundary is allowed to fall.

Prior work on lowering coroutines or \code{async} functions to hardware
exists and has not been surveyed for this paper. That is a gap, and it is
the first thing to close before Section~\ref{sec:e-next} is attempted,
because somebody may already have found where it breaks.

\section{What Comes Next}
\label{sec:e-next}

Nineteen probes prove that the design type-checks. Chisel is a working
compiler with a decade of use. The gap between those two facts is the
whole of what remains, and it should be closed in one place rather than
widened with more probes.

\subsection{The experiment}

Lower one function to a netlist.

\begin{lstlisting}[caption={The function to lower. Two operators, two awaits, one value live across a boundary.},label={lst:next-mac}]
async fn mac(a: U<32>, b: U<32>, prev: U<64>) -> U<64> {
    let p = mul(a, b).await;
    add(prev, p).await
}
\end{lstlisting}

The lowering has to produce, from that source and nothing else, a
two-stage pipeline in Verilog: a multiplier, a register holding \code{p}
across the boundary, an adder, and a valid signal that follows the data
through. It has to do so twice, once with \code{mul} bound to a
three-cycle implementation and once to a single-cycle one, and the second
netlist has to have one register fewer than the first without any edit to
Listing~\ref{lst:next-mac}.

That is the whole claim of Section~\ref{sec:e-pipeline} stated as a
test. If it passes, the behavioural half is real. If it cannot be made to
pass, the paper has rebuilt Chisel in Rust with better error messages,
which is worth having and is not what the project set out to do.

\subsection{What it needs}

A procedural macro on the function that reads its body as a syntax tree,
finds the awaits, and splits the body at each one into a stage. The
locals live across each split are the registers. The operators awaited
are instances whose latency is looked up in the configuration. The result
is a graph, and emitting Verilog from a graph is the part Chisel has
already shown to be routine.

Three things are known to be hard, and the experiment should be designed
to meet them rather than to avoid them.

\begin{itemize}
  \item \term{Control flow across an await.} A \code{when!} whose arms
        contain awaits is a pipeline whose stages are conditional. Calyx
        handles this with \code{if} in the control program; the lowering
        needs an equivalent, and it is the first place the async model
        might not fit.
  \item \term{Resource sharing.} Section~\ref{sec:e-pipeline} claims that
        concurrent invocations at the same await index share one
        operator. The lowering has to actually do that, or the claim is
        withdrawn.
  \item \term{Backpressure.} A stage that cannot proceed is an await that
        has not completed. The valid signal has to run backwards as well
        as forwards for that to be true in hardware, which is the elastic
        pipeline of Carloni~\cite{carloni2001}, and it has to be generated
        rather than assumed.
\end{itemize}

Since this section was first written the runtime has gained a
waveform: a probe registry named from the fields of units and a VCD
writer, so a design is watched in Surfer rather than in printed
traces. FST is the same probes with another writer, and the first
external crate the project would take. The same walk now writes a
structural netlist as well, a Verilog skeleton with modules,
registers, ports and wires and no behaviour; it sees fields only, so a
unit that takes its ports as parameters of \code{run} shows its
registers and nothing else. That gap is the argument for the
proc-macro route restated from the other side: the structure is
already recoverable at run time, and the behaviour is not.

\subsection{What it does not need}

A parser, a type checker, a module system, generics, a package manager or
an editor. That list is the argument for embedding, and it is why the
experiment is a macro of a few hundred lines rather than a compiler.
